\begin{table}[tbp]
\caption{Published plans and carbon-constrained portfolios in 2035. Values are total capacities in GW, including inherited assets. Carbon ranges span all four retained accepted deployment--architecture trajectories per system. Differences are carbon-constrained agent values minus the published plan, calculated before rounding. Brazil and South Africa use annual model budgets of 29 and 142~MtCO$_2$, respectively.}
\label{tab:expert-capacity}
\centering\small
\setlength{\tabcolsep}{5pt}
\begin{tabularx}{\textwidth}{@{}lXrrr@{}}
\toprule
System & Capacity & Expert plan & Carbon range & Difference \\
\midrule
Brazil & Wind & 49.55 & 120.19 & +70.64 \\
Brazil & Solar$^\dagger$ & 29.19 & 96.66 & +67.47 \\
Brazil & Battery power & 6.63 & 0.00 & -6.63 \\
South Africa & Wind & 26.39 & 5.57--6.24 & \([-20.82,-20.14]\) \\
South Africa & Solar$^\dagger$ & 21.16 & 63.02--64.24 & \([+41.86,+43.08]\) \\
\bottomrule
\end{tabularx}
\smallskip
\par\footnotesize $^\dagger$ Brazil's PDE 2035 centralized-solar category excludes 78.099~GW of micro-/mini-distributed generation. South Africa's IRP 2025 Solar PV column excludes 16.065~GW in its separate distributed-generation column, CSP and hybrid capacity. Model solar has no separate future distributed-generation rollout. The Brazilian values describe December 2035; the South African values combine the end-2025 base and scheduled additions through 2035. Demand, resources, inherited fleets, retirements, technology options and delivery schedules retain their respective study assumptions. Source locations and mappings are given in Methods and Appendix~\ref{app:regional-sources}.
\end{table}
